% =============================================================================
%  ملخص — Résumé en arabe (OPTIONNEL)
%
%  ⚠ Ce fichier ne compile PAS avec pdfLaTeX. Il exige XeLaTeX ou LuaLaTeX.
%
%  Si le projet Overleaf est déjà en XeLaTeX / LuaLaTeX, ajouter au préambule :
%
%      \usepackage{polyglossia}
%      \setmainlanguage{french}
%      \setotherlanguage{english}
%      \setotherlanguage{arabic}
%      \newfontfamily\arabicfont[Script=Arabic]{Amiri}   % ou Scheherazade New
%
%  ⚠ Le texte est dans un environnement `Arabic` AVEC UNE MAJUSCULE : en
%  minuscules, `\begin{arabic}` entre en collision avec la commande LaTeX
%  \arabic{compteur} et la compilation échoue sur « Missing number ». C'est
%  la forme documentée par polyglossia pour ce cas.
%
%  ⚠ Choisir une police arabe qui contient AUSSI la ponctuation latine : Geeza
%  Pro n'a pas le caractère « : » et laisse un rectangle vide. Amiri (TeX Live,
%  Overleaf) et Al Bayan (macOS) conviennent — vérifié.
%
%  puis, à l'endroit voulu du document :  % =============================================================================
%  ملخص — Résumé en arabe (OPTIONNEL)
%
%  ⚠ Ce fichier ne compile PAS avec pdfLaTeX. Il exige XeLaTeX ou LuaLaTeX.
%
%  Si le projet Overleaf est déjà en XeLaTeX / LuaLaTeX, ajouter au préambule :
%
%      \usepackage{polyglossia}
%      \setmainlanguage{french}
%      \setotherlanguage{english}
%      \setotherlanguage{arabic}
%      \newfontfamily\arabicfont[Script=Arabic]{Amiri}   % ou Scheherazade New
%
%  ⚠ Le texte est dans un environnement `Arabic` AVEC UNE MAJUSCULE : en
%  minuscules, `\begin{arabic}` entre en collision avec la commande LaTeX
%  \arabic{compteur} et la compilation échoue sur « Missing number ». C'est
%  la forme documentée par polyglossia pour ce cas.
%
%  ⚠ Choisir une police arabe qui contient AUSSI la ponctuation latine : Geeza
%  Pro n'a pas le caractère « : » et laisse un rectangle vide. Amiri (TeX Live,
%  Overleaf) et Al Bayan (macOS) conviennent — vérifié.
%
%  puis, à l'endroit voulu du document :  % =============================================================================
%  ملخص — Résumé en arabe (OPTIONNEL)
%
%  ⚠ Ce fichier ne compile PAS avec pdfLaTeX. Il exige XeLaTeX ou LuaLaTeX.
%
%  Si le projet Overleaf est déjà en XeLaTeX / LuaLaTeX, ajouter au préambule :
%
%      \usepackage{polyglossia}
%      \setmainlanguage{french}
%      \setotherlanguage{english}
%      \setotherlanguage{arabic}
%      \newfontfamily\arabicfont[Script=Arabic]{Amiri}   % ou Scheherazade New
%
%  ⚠ Le texte est dans un environnement `Arabic` AVEC UNE MAJUSCULE : en
%  minuscules, `\begin{arabic}` entre en collision avec la commande LaTeX
%  \arabic{compteur} et la compilation échoue sur « Missing number ». C'est
%  la forme documentée par polyglossia pour ce cas.
%
%  ⚠ Choisir une police arabe qui contient AUSSI la ponctuation latine : Geeza
%  Pro n'a pas le caractère « : » et laisse un rectangle vide. Amiri (TeX Live,
%  Overleaf) et Al Bayan (macOS) conviennent — vérifié.
%
%  puis, à l'endroit voulu du document :  % =============================================================================
%  ملخص — Résumé en arabe (OPTIONNEL)
%
%  ⚠ Ce fichier ne compile PAS avec pdfLaTeX. Il exige XeLaTeX ou LuaLaTeX.
%
%  Si le projet Overleaf est déjà en XeLaTeX / LuaLaTeX, ajouter au préambule :
%
%      \usepackage{polyglossia}
%      \setmainlanguage{french}
%      \setotherlanguage{english}
%      \setotherlanguage{arabic}
%      \newfontfamily\arabicfont[Script=Arabic]{Amiri}   % ou Scheherazade New
%
%  ⚠ Le texte est dans un environnement `Arabic` AVEC UNE MAJUSCULE : en
%  minuscules, `\begin{arabic}` entre en collision avec la commande LaTeX
%  \arabic{compteur} et la compilation échoue sur « Missing number ». C'est
%  la forme documentée par polyglossia pour ce cas.
%
%  ⚠ Choisir une police arabe qui contient AUSSI la ponctuation latine : Geeza
%  Pro n'a pas le caractère « : » et laisse un rectangle vide. Amiri (TeX Live,
%  Overleaf) et Al Bayan (macOS) conviennent — vérifié.
%
%  puis, à l'endroit voulu du document :  \input{frontmatter/resume_ar}
%
%  Si le projet reste en pdfLaTeX, NE PAS inclure ce fichier : le résumé
%  français et l'abstract anglais (frontmatter/resume.tex) suffisent, et un
%  résumé arabe mal rendu dessert le document plus qu'il ne le sert.
% =============================================================================

\chapter*{\textarabic{ملخص}}
\addcontentsline{toc}{chapter}{\textarabic{ملخص}}
\markboth{RÉSUMÉ (AR)}{RÉSUMÉ (AR)}

\begin{Arabic}

يتناول هذا المشروع، المنجَز عن بُعد لفائدة شركة \textenglish{EURAFRIC Information}، مسألة
توزيع رأس المال بين عدة أصول مالية: كيف يُوزَّع مبلغ من المال بين أربع شركات
مدرجة في بورصة الدار البيضاء وخمسة صناديق مؤشرات دولية، بحيث يُحصَّل أفضل عائد
ممكن مع التحكم في المخاطر.

تُجيب النظرية الكلاسيكية عن هذا السؤال منذ سنة 1952، غير أنها تفشل عملياً لأربعة
أسباب بنيوية: المقادير التي تقدّرها مشوَّشة، وسلوك الأسواق يتغيّر عبر الزمن،
ويختفي أثر التنويع في أوقات الأزمات بالذات، كما أن الأداء المقاس على بيانات
الماضي متفائل بشكل منهجي. يبني هذا العمل سلسلة معالجة كاملة تعالج هذه المشكلات
الأربع: ثلاث طبقات لتحضير البيانات والتحقق منها، ومحرك محاكاة زمنية تُثبَت
استحالة اطّلاعه على المستقبل باختبارات آلية لا بمجرد الادعاء، ثم أربع عائلات من
النماذج تُضاف واحدة تلو الأخرى: تقليص مصفوفة التغاير، والكشف غير المراقَب عن
أنظمة السوق بواسطة نموذج ماركوف المخفي، والتنبؤ بالعوائد بواسطة تجميعات
الأشجار، والتحسين الذي يأخذ تكاليف المعاملات بعين الاعتبار. ويُعامَل التقييم
بوصفه موضوعاً هندسياً قائماً بذاته: اختيار أمامي صارم (يسبق التدريبُ التحققَ
دائماً، مع تنقية وفترة حظر عند الحد)، وقطاع اختبار مجمَّد، ومقارنات
\textenglish{(paired)} مزدوجة بإعادة المعاينة بالكتل — أي اختبار الفرق بين
استراتيجيتين، وهو ما لا توفره الفترات الهامشية.

تُعرَض النتائج بحسب قوة الدليل. فقد تعرَّف كاشف الأنظمة على الأزمات الخمس
المدروسة جميعها، بأكثر من ثلاثة أضعاف معدله الأساسي ودون إنذارات دائمة؛ وتُعرض
هذه العلاقة بوصفها استكشافية لا تأكيدية لأن عدد الأزمات خمسة فقط. % Phrase remplacee INTEGRALEMENT (pas seulement le nombre) : l'ancienne
% affirmait une superiorite du systeme a regimes, sur donnees en devise mixte. La
% formulation ci-dessous est methodologiquement alignee sur les versions
% francaise et anglaise. A FAIRE RELIRE PAR UN ARABOPHONE AVANT REMISE.
بعد تصحيح العملة المرجعية باستخدام أسعار الصرف الرسمية لبنك المغرب، أصبح الفرق الملحوظ في نسبة شارب الصافية على محفظة \texttt{full\_2021} سالباً: ‎−10,47٪‎، إذ سجلت الاستراتيجية المشروطة بالنظام \texttt{regime\_conditional} نسبة ‎0,9571‎ مقابل ‎1,0690‎ لاستراتيجية \texttt{max\_sharpe}. هذا فرق وصفي ولا يثبت تفوقاً إحصائياً، لأنه لم يُجرَ اختبار مزدوج لهذا الفرق. أما عالم \texttt{etf\_2017} فيبقى مقوّماً بالدولار الأمريكي وغير متأثر بتصحيح العملة، مع فرق ملحوظ قدره ‎−1,6٪‎. وفي ظل الاختيار الأمامي
الصارم، احتوت جميع المقارنات المزدوجة الثماني على نافذة الاختبار المجمَّدة على
الصفر داخل فترات ثقتها: \textbf{لم تثبت أي أفضلية في الأداء}، لا لنماذج الإشارة
في مواجهة استراتيجية الأنظمة ولا في مواجهة التوزيع المتساوي. وهذا ليس دليلاً على
التكافؤ أيضاً، في غياب اختبار مخصص لذلك. كما استُكشفت خمسة مسارات مستقلة للتحسين
وأُغلقت، مما أثبت بوجه خاص أن دقة التنبؤ وأداء المحفظة هدفان متمايزان. وأخيراً
تبيّن أن سقف الوزن المفروض بوصفه قاعدة مهنية هو أنجع أداة تنظيم في النظام.
وتشمل المخرجات لوحة قيادة من صفحتين وواجهة برمجية \textenglish{REST}، يغطيها
أكثر من 580 اختباراً آلياً.

\vspace{0.6em}
\noindent\textbf{الكلمات المفتاحية:} تحسين المحفظة، التعلم الآلي، نموذج ماركوف
المخفي، مصفوفة التغاير، الاختبار الرجعي، التحقق خارج العينة، بورصة الدار
البيضاء، إدارة المخاطر.

\end{Arabic}

%
%  Si le projet reste en pdfLaTeX, NE PAS inclure ce fichier : le résumé
%  français et l'abstract anglais (frontmatter/resume.tex) suffisent, et un
%  résumé arabe mal rendu dessert le document plus qu'il ne le sert.
% =============================================================================

\chapter*{\textarabic{ملخص}}
\addcontentsline{toc}{chapter}{\textarabic{ملخص}}
\markboth{RÉSUMÉ (AR)}{RÉSUMÉ (AR)}

\begin{Arabic}

يتناول هذا المشروع، المنجَز عن بُعد لفائدة شركة \textenglish{EURAFRIC Information}، مسألة
توزيع رأس المال بين عدة أصول مالية: كيف يُوزَّع مبلغ من المال بين أربع شركات
مدرجة في بورصة الدار البيضاء وخمسة صناديق مؤشرات دولية، بحيث يُحصَّل أفضل عائد
ممكن مع التحكم في المخاطر.

تُجيب النظرية الكلاسيكية عن هذا السؤال منذ سنة 1952، غير أنها تفشل عملياً لأربعة
أسباب بنيوية: المقادير التي تقدّرها مشوَّشة، وسلوك الأسواق يتغيّر عبر الزمن،
ويختفي أثر التنويع في أوقات الأزمات بالذات، كما أن الأداء المقاس على بيانات
الماضي متفائل بشكل منهجي. يبني هذا العمل سلسلة معالجة كاملة تعالج هذه المشكلات
الأربع: ثلاث طبقات لتحضير البيانات والتحقق منها، ومحرك محاكاة زمنية تُثبَت
استحالة اطّلاعه على المستقبل باختبارات آلية لا بمجرد الادعاء، ثم أربع عائلات من
النماذج تُضاف واحدة تلو الأخرى: تقليص مصفوفة التغاير، والكشف غير المراقَب عن
أنظمة السوق بواسطة نموذج ماركوف المخفي، والتنبؤ بالعوائد بواسطة تجميعات
الأشجار، والتحسين الذي يأخذ تكاليف المعاملات بعين الاعتبار. ويُعامَل التقييم
بوصفه موضوعاً هندسياً قائماً بذاته: اختيار أمامي صارم (يسبق التدريبُ التحققَ
دائماً، مع تنقية وفترة حظر عند الحد)، وقطاع اختبار مجمَّد، ومقارنات
\textenglish{(paired)} مزدوجة بإعادة المعاينة بالكتل — أي اختبار الفرق بين
استراتيجيتين، وهو ما لا توفره الفترات الهامشية.

تُعرَض النتائج بحسب قوة الدليل. فقد تعرَّف كاشف الأنظمة على الأزمات الخمس
المدروسة جميعها، بأكثر من ثلاثة أضعاف معدله الأساسي ودون إنذارات دائمة؛ وتُعرض
هذه العلاقة بوصفها استكشافية لا تأكيدية لأن عدد الأزمات خمسة فقط. % Phrase remplacee INTEGRALEMENT (pas seulement le nombre) : l'ancienne
% affirmait une superiorite du systeme a regimes, sur donnees en devise mixte. La
% formulation ci-dessous est methodologiquement alignee sur les versions
% francaise et anglaise. A FAIRE RELIRE PAR UN ARABOPHONE AVANT REMISE.
بعد تصحيح العملة المرجعية باستخدام أسعار الصرف الرسمية لبنك المغرب، أصبح الفرق الملحوظ في نسبة شارب الصافية على محفظة \texttt{full\_2021} سالباً: ‎−10,47٪‎، إذ سجلت الاستراتيجية المشروطة بالنظام \texttt{regime\_conditional} نسبة ‎0,9571‎ مقابل ‎1,0690‎ لاستراتيجية \texttt{max\_sharpe}. هذا فرق وصفي ولا يثبت تفوقاً إحصائياً، لأنه لم يُجرَ اختبار مزدوج لهذا الفرق. أما عالم \texttt{etf\_2017} فيبقى مقوّماً بالدولار الأمريكي وغير متأثر بتصحيح العملة، مع فرق ملحوظ قدره ‎−1,6٪‎. وفي ظل الاختيار الأمامي
الصارم، احتوت جميع المقارنات المزدوجة الثماني على نافذة الاختبار المجمَّدة على
الصفر داخل فترات ثقتها: \textbf{لم تثبت أي أفضلية في الأداء}، لا لنماذج الإشارة
في مواجهة استراتيجية الأنظمة ولا في مواجهة التوزيع المتساوي. وهذا ليس دليلاً على
التكافؤ أيضاً، في غياب اختبار مخصص لذلك. كما استُكشفت خمسة مسارات مستقلة للتحسين
وأُغلقت، مما أثبت بوجه خاص أن دقة التنبؤ وأداء المحفظة هدفان متمايزان. وأخيراً
تبيّن أن سقف الوزن المفروض بوصفه قاعدة مهنية هو أنجع أداة تنظيم في النظام.
وتشمل المخرجات لوحة قيادة من صفحتين وواجهة برمجية \textenglish{REST}، يغطيها
أكثر من 580 اختباراً آلياً.

\vspace{0.6em}
\noindent\textbf{الكلمات المفتاحية:} تحسين المحفظة، التعلم الآلي، نموذج ماركوف
المخفي، مصفوفة التغاير، الاختبار الرجعي، التحقق خارج العينة، بورصة الدار
البيضاء، إدارة المخاطر.

\end{Arabic}

%
%  Si le projet reste en pdfLaTeX, NE PAS inclure ce fichier : le résumé
%  français et l'abstract anglais (frontmatter/resume.tex) suffisent, et un
%  résumé arabe mal rendu dessert le document plus qu'il ne le sert.
% =============================================================================

\chapter*{\textarabic{ملخص}}
\addcontentsline{toc}{chapter}{\textarabic{ملخص}}
\markboth{RÉSUMÉ (AR)}{RÉSUMÉ (AR)}

\begin{Arabic}

يتناول هذا المشروع، المنجَز عن بُعد لفائدة شركة \textenglish{EURAFRIC Information}، مسألة
توزيع رأس المال بين عدة أصول مالية: كيف يُوزَّع مبلغ من المال بين أربع شركات
مدرجة في بورصة الدار البيضاء وخمسة صناديق مؤشرات دولية، بحيث يُحصَّل أفضل عائد
ممكن مع التحكم في المخاطر.

تُجيب النظرية الكلاسيكية عن هذا السؤال منذ سنة 1952، غير أنها تفشل عملياً لأربعة
أسباب بنيوية: المقادير التي تقدّرها مشوَّشة، وسلوك الأسواق يتغيّر عبر الزمن،
ويختفي أثر التنويع في أوقات الأزمات بالذات، كما أن الأداء المقاس على بيانات
الماضي متفائل بشكل منهجي. يبني هذا العمل سلسلة معالجة كاملة تعالج هذه المشكلات
الأربع: ثلاث طبقات لتحضير البيانات والتحقق منها، ومحرك محاكاة زمنية تُثبَت
استحالة اطّلاعه على المستقبل باختبارات آلية لا بمجرد الادعاء، ثم أربع عائلات من
النماذج تُضاف واحدة تلو الأخرى: تقليص مصفوفة التغاير، والكشف غير المراقَب عن
أنظمة السوق بواسطة نموذج ماركوف المخفي، والتنبؤ بالعوائد بواسطة تجميعات
الأشجار، والتحسين الذي يأخذ تكاليف المعاملات بعين الاعتبار. ويُعامَل التقييم
بوصفه موضوعاً هندسياً قائماً بذاته: اختيار أمامي صارم (يسبق التدريبُ التحققَ
دائماً، مع تنقية وفترة حظر عند الحد)، وقطاع اختبار مجمَّد، ومقارنات
\textenglish{(paired)} مزدوجة بإعادة المعاينة بالكتل — أي اختبار الفرق بين
استراتيجيتين، وهو ما لا توفره الفترات الهامشية.

تُعرَض النتائج بحسب قوة الدليل. فقد تعرَّف كاشف الأنظمة على الأزمات الخمس
المدروسة جميعها، بأكثر من ثلاثة أضعاف معدله الأساسي ودون إنذارات دائمة؛ وتُعرض
هذه العلاقة بوصفها استكشافية لا تأكيدية لأن عدد الأزمات خمسة فقط. % Phrase remplacee INTEGRALEMENT (pas seulement le nombre) : l'ancienne
% affirmait une superiorite du systeme a regimes, sur donnees en devise mixte. La
% formulation ci-dessous est methodologiquement alignee sur les versions
% francaise et anglaise. A FAIRE RELIRE PAR UN ARABOPHONE AVANT REMISE.
بعد تصحيح العملة المرجعية باستخدام أسعار الصرف الرسمية لبنك المغرب، أصبح الفرق الملحوظ في نسبة شارب الصافية على محفظة \texttt{full\_2021} سالباً: ‎−10,47٪‎، إذ سجلت الاستراتيجية المشروطة بالنظام \texttt{regime\_conditional} نسبة ‎0,9571‎ مقابل ‎1,0690‎ لاستراتيجية \texttt{max\_sharpe}. هذا فرق وصفي ولا يثبت تفوقاً إحصائياً، لأنه لم يُجرَ اختبار مزدوج لهذا الفرق. أما عالم \texttt{etf\_2017} فيبقى مقوّماً بالدولار الأمريكي وغير متأثر بتصحيح العملة، مع فرق ملحوظ قدره ‎−1,6٪‎. وفي ظل الاختيار الأمامي
الصارم، احتوت جميع المقارنات المزدوجة الثماني على نافذة الاختبار المجمَّدة على
الصفر داخل فترات ثقتها: \textbf{لم تثبت أي أفضلية في الأداء}، لا لنماذج الإشارة
في مواجهة استراتيجية الأنظمة ولا في مواجهة التوزيع المتساوي. وهذا ليس دليلاً على
التكافؤ أيضاً، في غياب اختبار مخصص لذلك. كما استُكشفت خمسة مسارات مستقلة للتحسين
وأُغلقت، مما أثبت بوجه خاص أن دقة التنبؤ وأداء المحفظة هدفان متمايزان. وأخيراً
تبيّن أن سقف الوزن المفروض بوصفه قاعدة مهنية هو أنجع أداة تنظيم في النظام.
وتشمل المخرجات لوحة قيادة من صفحتين وواجهة برمجية \textenglish{REST}، يغطيها
أكثر من 580 اختباراً آلياً.

\vspace{0.6em}
\noindent\textbf{الكلمات المفتاحية:} تحسين المحفظة، التعلم الآلي، نموذج ماركوف
المخفي، مصفوفة التغاير، الاختبار الرجعي، التحقق خارج العينة، بورصة الدار
البيضاء، إدارة المخاطر.

\end{Arabic}

%
%  Si le projet reste en pdfLaTeX, NE PAS inclure ce fichier : le résumé
%  français et l'abstract anglais (frontmatter/resume.tex) suffisent, et un
%  résumé arabe mal rendu dessert le document plus qu'il ne le sert.
% =============================================================================

\chapter*{\textarabic{ملخص}}
\addcontentsline{toc}{chapter}{\textarabic{ملخص}}
\markboth{RÉSUMÉ (AR)}{RÉSUMÉ (AR)}

\begin{Arabic}

يتناول هذا المشروع، المنجَز عن بُعد لفائدة شركة \textenglish{EURAFRIC Information}، مسألة
توزيع رأس المال بين عدة أصول مالية: كيف يُوزَّع مبلغ من المال بين أربع شركات
مدرجة في بورصة الدار البيضاء وخمسة صناديق مؤشرات دولية، بحيث يُحصَّل أفضل عائد
ممكن مع التحكم في المخاطر.

تُجيب النظرية الكلاسيكية عن هذا السؤال منذ سنة 1952، غير أنها تفشل عملياً لأربعة
أسباب بنيوية: المقادير التي تقدّرها مشوَّشة، وسلوك الأسواق يتغيّر عبر الزمن،
ويختفي أثر التنويع في أوقات الأزمات بالذات، كما أن الأداء المقاس على بيانات
الماضي متفائل بشكل منهجي. يبني هذا العمل سلسلة معالجة كاملة تعالج هذه المشكلات
الأربع: ثلاث طبقات لتحضير البيانات والتحقق منها، ومحرك محاكاة زمنية تُثبَت
استحالة اطّلاعه على المستقبل باختبارات آلية لا بمجرد الادعاء، ثم أربع عائلات من
النماذج تُضاف واحدة تلو الأخرى: تقليص مصفوفة التغاير، والكشف غير المراقَب عن
أنظمة السوق بواسطة نموذج ماركوف المخفي، والتنبؤ بالعوائد بواسطة تجميعات
الأشجار، والتحسين الذي يأخذ تكاليف المعاملات بعين الاعتبار. ويُعامَل التقييم
بوصفه موضوعاً هندسياً قائماً بذاته: اختيار أمامي صارم (يسبق التدريبُ التحققَ
دائماً، مع تنقية وفترة حظر عند الحد)، وقطاع اختبار مجمَّد، ومقارنات
\textenglish{(paired)} مزدوجة بإعادة المعاينة بالكتل — أي اختبار الفرق بين
استراتيجيتين، وهو ما لا توفره الفترات الهامشية.

تُعرَض النتائج بحسب قوة الدليل. فقد تعرَّف كاشف الأنظمة على الأزمات الخمس
المدروسة جميعها، بأكثر من ثلاثة أضعاف معدله الأساسي ودون إنذارات دائمة؛ وتُعرض
هذه العلاقة بوصفها استكشافية لا تأكيدية لأن عدد الأزمات خمسة فقط. % Phrase remplacee INTEGRALEMENT (pas seulement le nombre) : l'ancienne
% affirmait une superiorite du systeme a regimes, sur donnees en devise mixte. La
% formulation ci-dessous est methodologiquement alignee sur les versions
% francaise et anglaise. A FAIRE RELIRE PAR UN ARABOPHONE AVANT REMISE.
بعد تصحيح العملة المرجعية باستخدام أسعار الصرف الرسمية لبنك المغرب، أصبح الفرق الملحوظ في نسبة شارب الصافية على محفظة \texttt{full\_2021} سالباً: ‎−10,47٪‎، إذ سجلت الاستراتيجية المشروطة بالنظام \texttt{regime\_conditional} نسبة ‎0,9571‎ مقابل ‎1,0690‎ لاستراتيجية \texttt{max\_sharpe}. هذا فرق وصفي ولا يثبت تفوقاً إحصائياً، لأنه لم يُجرَ اختبار مزدوج لهذا الفرق. أما عالم \texttt{etf\_2017} فيبقى مقوّماً بالدولار الأمريكي وغير متأثر بتصحيح العملة، مع فرق ملحوظ قدره ‎−1,6٪‎. وفي ظل الاختيار الأمامي
الصارم، احتوت جميع المقارنات المزدوجة الثماني على نافذة الاختبار المجمَّدة على
الصفر داخل فترات ثقتها: \textbf{لم تثبت أي أفضلية في الأداء}، لا لنماذج الإشارة
في مواجهة استراتيجية الأنظمة ولا في مواجهة التوزيع المتساوي. وهذا ليس دليلاً على
التكافؤ أيضاً، في غياب اختبار مخصص لذلك. كما استُكشفت خمسة مسارات مستقلة للتحسين
وأُغلقت، مما أثبت بوجه خاص أن دقة التنبؤ وأداء المحفظة هدفان متمايزان. وأخيراً
تبيّن أن سقف الوزن المفروض بوصفه قاعدة مهنية هو أنجع أداة تنظيم في النظام.
وتشمل المخرجات لوحة قيادة من صفحتين وواجهة برمجية \textenglish{REST}، يغطيها
أكثر من 580 اختباراً آلياً.

\vspace{0.6em}
\noindent\textbf{الكلمات المفتاحية:} تحسين المحفظة، التعلم الآلي، نموذج ماركوف
المخفي، مصفوفة التغاير، الاختبار الرجعي، التحقق خارج العينة، بورصة الدار
البيضاء، إدارة المخاطر.

\end{Arabic}
